Ein Punkt am Äquator legt während einer Erdumdrehung ($T = \TeO$) den Erdumfang $2\pi R$ zurück:
\begin{align}
 v &= \veF \\
   &= \frac{2\pi\cdot\rE}{\Te} \\
   &= \ve \approx \result{\veP{0}{3}}
\end{align}
Genau genommen dauert eine Umdrehung der Erde gegenüber den Fixsternen einen Sterntag, $T_\text{s} = \TsO$ (FoTaBe S.~205). Damit erhält man $\vsdP{0}{3}$; der Unterschied ist vernachlässigbar.
